\section{Ergebnisse}
In diesem Kapitel werden die zuvor vorgestellten C2-Frameworks anhand verschiedener Kriterien miteinander verglichen. Ziel ist es, die Stärken und Schwächen der einzelnen Tools herauszuarbeiten und eine fundierte Bewertung ihrer Eignung für unterschiedliche Anwendungsfälle im Red Teaming zu ermöglichen. Die Vergleichskriterien umfassen unter anderem die Benutzerfreundlichkeit, die Flexibilität bei der Anpassung von Payloads, die Fähigkeit zur Umgehung von Sicherheitslösungen sowie die unterstützten Betriebssysteme und Protokolle. Durch diese vergleichende Analyse soll ein umfassender Überblick über die Leistungsfähigkeit der verschiedenen C2-Frameworks gegeben werden, um eine informierte Entscheidung bei der Auswahl des geeigneten Tools für spezifische Angriffsszenarien zu erleichtern.
\subsection{Techischer Vergleich}
- Kommunikationsprotokolle
- Payload-Generierung
- Modularität
- Erweiterbarkeit
\subsection{Detection und Blue-Team-Sicht}
- Signaturen
- Verhaltensbasierte Erkennung
- Netzwerkverkehrsanalyse
\subsection{Einsatzszenarien}
- Wann Sliver sinnvoller ist
- Wann msf ausreicht
- warum cobalt strike in manchen Fällen die bessere Wahl sein könnte